%15 min preso!
\documentclass[xcolor=table,aspectratio=169]{beamer}
\usepackage{beamerthemesplit}
\usepackage{wrapfig}
\usetheme{SPbGU}
\usepackage{pdfpages}
\usepackage{amsmath}
\usepackage{cmap}
\usepackage[T2A]{fontenc}
\usepackage[utf8]{inputenc}
\usepackage[english]{babel}
\usepackage{indentfirst}
\usepackage{amsmath}
\usepackage{tikz}
\usepackage{multirow}
\usepackage[noend]{algpseudocode}
\usepackage{algorithm}
\usepackage{algorithmicx}
\usepackage{fancyvrb}
\usepackage{hyperref} 
\usetikzlibrary{calc}
\usetikzlibrary{shapes, backgrounds}
\usetikzlibrary{arrows,automata}
\usetikzlibrary{positioning}
\usetikzlibrary{fit}
\usetikzlibrary{shapes.callouts}
\usetikzlibrary{shapes.misc}
\usepackage{xparse}
\usepackage{fontawesome}

\usepackage{etoolbox,refcount}
\usepackage{multicol}

\usepackage{tabularx}
\newcolumntype{Y}{>{\raggedleft\arraybackslash}X}

\renewcommand{\thealgorithm}{}

\newtheorem{mytheorem}{Theorem}
\renewcommand{\thealgorithm}{}

\newcommand{\tikzmark}[1]{\tikz[overlay,remember picture] \node (#1) {};}
\def\Put(#1,#2)#3{\leavevmode\makebox(0,0){\put(#1,#2){#3}}}

\newcommand{\ltz}{$< 1$}

\tikzset{
    state/.style={
           rectangle,
           rounded corners,
           draw=black, very thick,
           minimum height=2em,
           inner sep=2pt,
           text centered,
           },
}

\tikzset{
    invisible/.style={opacity=0,text opacity=0},
    visible on/.style={alt=#1{}{invisible}},
    alt/.code args={<#1>#2#3}{%
      \alt<#1>{\pgfkeysalso{#2}}{\pgfkeysalso{#3}} % \pgfkeysalso doesn't change the path
    },
}

\tikzset{cross/.style={cross out, draw=black, minimum size=2*(#1-\pgflinewidth), inner sep=0pt, outer sep=0pt, ultra thick},
%default radius will be 1pt. 
cross/.default={1pt}}

\NewDocumentCommand{\mycallout}{r<> O{opacity=0.8,text opacity=1} m m m}{%
\tikz[remember picture, overlay]\node[align=center, fill=cyan!20, text width=#5cm,
#2,visible on=<#1>, rounded corners,
draw,rectangle callout,anchor=pointer,callout relative pointer={(290:0.5cm)}]
at (#3) {#4};
}

\NewDocumentCommand{\mycalloutR}{r<> O{opacity=0.8,text opacity=1} m m m}{%
\tikz[remember picture, overlay]\node[align=center, fill=cyan!20, text width=#5cm,
#2,visible on=<#1>, rounded corners,
draw,rectangle callout,anchor=pointer,callout relative pointer={(30:0.8cm)}]
at (#3) {#4};
}


%callout relative pointer={(230:0.5cm)}]

\newcounter{countitems}
\newcounter{nextitemizecount}
\newcommand{\setupcountitems}{%
  \stepcounter{nextitemizecount}%
  \setcounter{countitems}{0}%
  \preto\item{\stepcounter{countitems}}%
}
\makeatletter
\newcommand{\computecountitems}{%
  \edef\@currentlabel{\number\c@countitems}%
  \label{countitems@\number\numexpr\value{nextitemizecount}-1\relax}%
}
\newcommand{\nextitemizecount}{%
  \getrefnumber{countitems@\number\c@nextitemizecount}%
}
\newcommand{\previtemizecount}{%
  \getrefnumber{countitems@\number\numexpr\value{nextitemizecount}-1\relax}%
}
\makeatother    
\newenvironment{AutoMultiColItemize}{%
\ifnumcomp{\nextitemizecount}{>}{3}{\begin{multicols}{2}}{}%
\setupcountitems\begin{itemize}}%
{\end{itemize}%
\unskip\computecountitems\ifnumcomp{\previtemizecount}{>}{3}{\end{multicols}}{}}


\beamertemplatenavigationsymbolsempty

\title[Разреженная линейная алгебра: зачем и как]{Линейная алгебра как путь к высокопроизводительному анализу данных: программные и аппаратные средства, их проблемы и возможности платформы RISC-V}
\institute[СПбГУ]{
Санкт-Петербургский Государственный Университет
}

% То, что в квадратных скобках, отображается в левом нижнем углу.
\author[Семён Григорьев]{Семён Григорьев}

\date{20 сентября 2024г.}


%Я предлагаю сначала в общем рассказать интересы и компетенции группы, 
%что научная группа сделала за 3 месяца, 
%потом потратить 1 страницу презентации на каждого из студентов (что сделал, почему важно для Yiming, какие планы). 
%В конце перейти к планам на год.

\begin{document}
{
\begin{frame}[fragile]
  \begin{table}
  \centering
  %\includegraphics[height=1.5cm]{pictures/SPbGU_Logo.png}
  \begin{tabularx}{\linewidth}{XcX}
    \includegraphics[height=1.6cm]{pictures/YRS.pdf} \hfill
    & 
    & \hfill \includegraphics[height=1.6cm]{pictures/SPbGU_Logo.png}
  \end{tabularx}
  \end{table}
  \titlepage
\end{frame}
}

\begin{frame}[fragile]
  \frametitle{Семён Григорьев}
  \begin{itemize}    
    \item Область интересов: алгоритмы анализа графов, высокопроизводительная разреженная линейная алгебра, теория формальных языков и её применение для анализа данных
    \vspace{2cm}
    \item Доцент кафедры системного программирования СПбГУ
    \item Почта: rsdpisuy@gmail.com 
    \item ТГ: @rsdpisuy
  \end{itemize}
\end{frame}

\begin{frame}[fragile]
  \frametitle{Разреженная линейная алгебра}
  \begin{itemize}
    \item Линейная алгебра: матрицы, вектора, операции над ними
    \begin{itemize}
      \item Операции естественным образом распараллеливаются по данным: эффективные реализации для многоядерных систем, GPGPU, и т.д.
      \item Абстракция по операциям: (полу)кольца, моноиды, $\ldots$
    \end{itemize}
    \item \textbf{Разреженная} линейная алгебра: в матрице или векторе много одинаковых элементов
    \begin{itemize}
      \item Часто говорят что в матрице (векторе) много <<нейтральных элементов>>, <<нулей>> или что-то подобное, но это не всегда так
    \end{itemize}
    \item Хотим не хранить одинаковые элементы
    \begin{itemize}
      \item Специальные структуры для хранения матриц и векторов\footnote{COO, CSR, CSC, DCSR, Quad-Tree, $\ldots$}
      \item Специальные алгоритмы для выполнения операций\footnote{Не забываем про параллельность}
    \end{itemize}
  \end{itemize}
\end{frame}


\begin{frame}[fragile]
  \frametitle{Области применения}
  \begin{minipage}{0.46\textwidth}
    \begin{itemize}
    \item Машинное обучение
    \begin{itemize}
      \item Разреженное внимание (sparse attention)
      \item Графовые нейронные сети
    \end{itemize}
    \item Робототехника
    \begin{itemize}
      \item Задачи навигации
      \item $\ldots$
    \end{itemize}
    \item Численные методы
    \begin{itemize}
      \item Разреженные системы уравнений
      \item $\ldots$
    \end{itemize}
    \item $\ldots$
  \end{itemize}
  \end{minipage}
  \begin{minipage}{0.46\textwidth}
    \begin{itemize}
      \item Анализ графов
      \begin{itemize}
        \item Графовые базы данных
        \item Анализ социальных, банковских и других сетей
        \item Статический анализ кода
        \item Биоинформатика
      \end{itemize}
    \end{itemize}
  \end{minipage}
  
\end{frame}


\begin{frame}[fragile]
  \frametitle{GraphBLAS}
  Стандарт C API, определяющий набор примитивов обобщённой разреженной линейной алгебры и операций над ними  
  \begin{itemize}
    \item Домашняя страница: \href{https://graphblas.org/}{https://graphblas.org/} 
    \item Коллекция полезных ссылок по теме (GraphBLAS-Pointers): \href{https://github.com/GraphBLAS/GraphBLAS-Pointers}{https://github.com/GraphBLAS/GraphBLAS-Pointers}  
    \item Эталонная реализация на C (SuiteSparse:GraphBLAS): \href{https://github.com/DrTimothyAldenDavis/GraphBLAS}{https://github.com/DrTimothyAldenDavis/GraphBLAS} 
    \item Обёртка для Python (python-graphblas): \href{https://github.com/python-graphblas/python-graphblas}{https://github.com/python-graphblas/python-graphblas}
    \item Для распределённой обработки (CombBLAS): \href{https://github.com/PASSIONLab/CombBLAS}{https://github.com/PASSIONLab/CombBLAS}
    \item Графовая база данных (FalkorDB): \href{https://github.com/FalkorDB/falkordb}{https://github.com/FalkorDB/falkordb}
  \end{itemize}
\end{frame}

\begin{frame}[fragile]
  \frametitle{Специализированные решения для разреженной линейной алгебры}
  \begin{itemize}
    \item Dedicated Hardware Accelerators for Processing of Sparse Matrices and Vectors: A Survey, 2024 год
    \item A Survey of Accelerating Parallel Sparse Linear Algebra, 2023 год
    \item A Systematic Literature Survey of Sparse Matrix-Vector Multiplication, 2024 год
    \item GraphLily (подмножество GraphBLAS на FPGA): \href{https://github.com/cornell-zhang/GraphLily}{https://github.com/cornell-zhang/GraphLily}
  \end{itemize}
\end{frame}

\begin{frame}[fragile]
  \frametitle{А что про RISC-V?}
  \begin{itemize}    
    \item Идёт работа над расширениями\footnote{Оставим в покое RVV, Integrated Matrix Extension, XuanTie Matrix Extension}
    \begin{itemize}
      \item IndexMAC: A Custom RISC-V Vector Instruction to Accelerate Structured-Sparse Matrix Multiplications, 2024 год
      \item Optimizations for Very Long and Sparse Vector Operations on a RISC-V VPU: A Work-in-Progress, 2023 год
    \end{itemize}
    \item Vortex
    \begin{itemize}
      \item RISC-V GPGPU: \href{https://github.com/vortexgpgpu/vortex}{https://github.com/vortexgpgpu/vortex}
      \item Молодой, перспективный, $\ldots$
    \end{itemize}
  \end{itemize}
\end{frame}

\begin{frame}[fragile]
  \frametitle{Заключение}  
  \begin{itemize}
    \item Высокопроизводительная разреженная линейная алгебра $\Rightarrow$ высокопроизводительные приложения
    \begin{itemize}
      \item Машинное обучение
      \item Графовые базы данных
      \item Анализ социальных, банковских и других сетей
      \item Анализ кода
      \item $\ldots$
    \end{itemize}
    \item Сделать \textbf{разреженную} линейную алгебру высокопроизводительной \textbf{сложно}
    \begin{itemize}
      \item Нерегулярный доступ к данным
      \item Хорошая алгебра --- обобщённая алгебра
      \item Сложности с балансировкой нагрузки
      \item $\ldots$
    \end{itemize}
    \item Но люди пытаются
    \begin{itemize}
      \item Даже институты для этого создают\footnote{Sparsitute: A mathematical Institute for Sparse Computations in Science and Engineering, \href{https://sparsitute.lbl.gov/}{https://sparsitute.lbl.gov/}}
    \end{itemize}
  \end{itemize}
\end{frame}


\end{document}
